%% capitoli/cap6_unitarieta.tex
%% In testa al file, prima di scrivere, incollare l'estratto di
%% output/F4-indice.md relativo a questo capitolo (tesi sostenuta,
%% fonti, lunghezza stimata): serve a controllare gli scostamenti.
%% Convenzione delle etichette: sez:cap6:<nome>, tab:cap6:<nome>.

% Capitolo della presa di posizione: poggia sugli indici ricostruiti nei
% capitoli IV e V. Non si scrive prima che quelli siano chiusi.
% La matrice degli indici va in tabella: \label{tab:cap6:indici}

\chapter{Funzione unitaria o pluralità di funzioni consultive}
\label{cap:cap6}

\section{La tesi unitaria: argomenti e punti di forza}
\label{sez:cap6:unitaria}


\section{La tesi pluralista: eterogeneità di legittimazione, oggetto ed effetti}
\label{sez:cap6:pluralista}


\section{La via intermedia: un genus con specie a effetti differenziati}
\label{sez:cap6:intermedia}


\section{Presa di posizione e verifica sugli indici}
\label{sez:cap6:posizione}


\section{Il rischio della co-amministrazione e il problema dell'imparzialità}
\label{sez:cap6:coamministrazione}


